%!TeX program = lualatex
\renewcommand{\thelektion}{L07 \textperiodcentered{} Dateiformate \& Speicherplatz}

\begin{frame}\lessontitle{Lektion 07}{Dateiformate \& Speicherplatz}{Wie man aus Bits ein Dateiformat erkennt \textperiodcentered{} Skript, Kapitel 3}\end{frame}

\begin{frame}\FDUtitle{Eine einfache Frage}
\begin{center}
Wir wissen: Computer speichern nur Nullen und Einsen.
\vspace{4mm}

Woher \enquote{weiss} ein Computer, ob eine Bitfolge ein Bild, eine Zahl, ein Buchstabe oder etwas anderes ist?
\end{center}
\end{frame}

\begin{frame}\FDUtitle{Dateiformate geben Bedeutung}
\begin{table}[H]
	\centering
	\small
	\begin{tabular}{l|p{6.3cm}}
		\textbf{Datei} & \textbf{Bedeutung (Beispiel)} \\\midrule
		\texttt{.txt} & Zeichen: $<|\ldots$ \\
		\texttt{.doc} & Formatierbefehl: \enquote{Einzug 1.75\,cm} \\
		\texttt{.xls} & Zahl 9662249 in einer Zelle \\
		\texttt{.gif} & Bildpunkt, indexiert (256 Farben, 1 Byte/Pixel) \\
		\texttt{.bmp} & Bildpunkt, direkt (16.7\,Mio.\ Farben, 3 Bytes/Pixel) \\
	\end{tabular}
\end{table}
\begin{center}Dieselbe Bitfolge kann je nach Dateiformat ganz unterschiedlich interpretiert werden!\end{center}
\end{frame}

\begin{frame}\FDUtitle{Speichergrössen}
\begin{table}[H]
	\footnotesize
	\centering
	\begin{tabular}{@{}l|rr|l@{}}
		\toprule
		\textbf{Masseinheit} & \multicolumn{2}{l|}{\textbf{Dezimalsystem}} & Grössenordnung                                   \\ \midrule
		KB (Kilobyte)        & $10^3$ B                               & 1'000 B             & eine Text-Datei            \\
		MB (Megabyte)        & $10^6$ B                                    & 1'000'000 B         & eine Musik-Datei           \\
		GB (Gigabyte)        & $10^9$ B                                    & 1'000'000'000 B     & eine Video-Datei           \\
		TB (Terabyte)        & $10^{12}$ B                                 & 1'000'000'000'000 B & kleiner Firmen-Server      \\
		PB (Petabyte)        & $10^{15}$ B                                 & ...                 & Facebook-Server (Meta)     \\
		EB (Exabyte)         & $10^{18}$ B                                 & ...                 & alle CERN-Daten            \\
		ZB (Zettabyte)       & $10^{21}$ B                                 & ...                 & alle Daten ($\sim$ 100 ZB) \\ \bottomrule
	\end{tabular}
\end{table}
\only<2->{
\begin{center}
	Wie nennt man ein \enquote{Megagramm} geläufiger? $\rightarrow$ eine Tonne
\end{center}}
\end{frame}

\begin{frame}\FDUtitle{Diskussion}
\begin{center}
Weshalb haben Festplatten häufig eine leicht kleinere Speichergrösse als angegeben, insbesondere unter Windows?
\end{center}
\only<2->{
\begin{center}\small Hersteller rechnen mit $1\,\mathrm{GB}=10^9$ Byte, das Betriebssystem oft mit $1\,\mathrm{GiB}=2^{30}$ Byte -- ein kleiner, aber spürbarer Unterschied.\end{center}
}
\end{frame}

\begin{frame}{Auftrag}{Kapitel 3: Dateiformate \& Speicherbedarf}
\aufgabebox{\textbf{Aufgabe 3.15}: Video anschauen (Passwort: BinCode)}\\[3mm]
\aufgabebox{\textbf{Aufgabe 3.16}: Üblicher/gebräuchlicher Name für ein \enquote{Megagramm}}\\[3mm]
\challengebox{\textbf{Aufgabe 3.17 (Challenge)}: Warum Festplatten unter Windows kleiner wirken (Apple-Artikel lesen)}
\end{frame}
